\section{Introduction}
% Useful links:
% \begin{itemize}
% 	\item Online Table Generator \\ \url{http://www.tablesgenerator.com/latex_tables}
% 	\item Flowcharts and drawings easily online \\ \url{https://www.draw.io}
% 	\item Minipage Help \\ \url{http://www.namsu.de/Extra/befehle/Minipage.html}
% 	\item Online Formula Editor \\ \url{https://www.matheretter.de/tools/formeleditor/}
% 	\item Online Formula Editor with PDF Export \\ \url{https://www.latex4technics.com}
% \end{itemize}

% % Listing of vectors



% % QR Code

% For linking a video, a QR code is helpful. Here it is shown as an example.
% \marginpar{
% 	\qrset{link,height=1.5cm}
% 	\qrcode{https://www.youtube.com/watch?v=fox6qA6SuYc}}

% Table with diagonal line


Modern industrial environments demand flexible, efficient, and adaptive automation systems. As production cycles shorten and product variants grow, traditional robot programming becomes a bottleneck, requiring specialized knowledge and significant time investment \cite{calinon2009robot}. In response, Human--Robot Collaboration (HRC) and, in particular, Programming by Demonstration (PbD) have emerged as promising alternatives \cite{calinon2009robot}.

PbD enables non-expert users to teach robots new tasks through physical or visual demonstration, bypassing manual coding. Its effectiveness, however, depends on accurate perception and replication of human-guided interactions in real-world settings \cite{calinon2009robot}.

Augmented Reality (AR) offers new ways to enhance PbD by enabling spatially aware, intuitive interaction via gaze, gestures, or other natural inputs \cite{lotsaris2021ar}. Yet, a major technical challenge remains: reliable, low-latency six-degree-of-freedom (6D) pose estimation of physical objects. Industrial conditions such as occlusion, dynamic lighting, and object variability still hinder robustness and accuracy, while many existing solutions lack the responsiveness required for interactive operator guidance \cite{thalhammer2024challenges}.

Recent advances in artificial intelligence (AI)-based pose estimation and the arrival of consumer-grade AR hardware like the Apple Vision Pro create an opportunity to address these issues. By combining deep learning with modern AR spatial tracking, it is now feasible to build user-friendly, interactive systems for robot instruction \cite{lotsaris2021ar,wen2024foundationpose}.

This thesis develops a proof-of-concept AR application for intuitive robot programming via low-latency 6D pose estimation. Using the Apple Vision Pro, the system supports gaze- and gesture-based interaction, object or region selection, backend pose estimation, and AR visualization. Continuous sensor data streaming targets interactive update rates, with evaluation focusing on accuracy, usability, and responsiveness.

The remainder of this thesis is structured as follows. Section~2 reviews the state of the art in PbD, AR-based robot programming, and 6D pose estimation. Section~3 defines the overall objective of the work and outlines the methodological approach. Section~4 presents the system design, including the backend pose estimation pipeline and prototyping strategy. Section~5 describes the implementation of the mixed reality application on visionOS. Section~6 discusses the evaluation methodology and preliminary results. Section~7 concludes the thesis and provides an outlook on future research directions and potential extensions of the system.
